\section{Complex tori}
\label{mumford-I1}
\dcref{M:ch1:I.1}

\begin{convention}[Analytic standing hypotheses]
\label{mumford-conv-complex-analytic}
\dcref{M:ch1:I.1}
Throughout this chapter $V$ is a complex vector space of dimension $g$ and
$\Lambda\subset V$ is a discrete subgroup spanning $V$ over $\C$.  A complex
torus is the quotient $V/\Lambda$ with its quotient complex structure.  A
compact connected complex Lie group is always understood to have holomorphic
multiplication and inverse.
\notready
\end{convention}

\begin{lemma}[Compactness of a lattice quotient]
\label{mumford-lem-compact-lattice-quotient}
\dcref{M:ch1:I.1}
If $\Lambda$ is a lattice in $V$, then $V/\Lambda$ is a compact connected
complex manifold of dimension $g$, and the quotient map is a holomorphic
covering with deck group $\Lambda$.
\notready
\end{lemma}

\begin{lemma}[Compact holomorphic functions]
\label{mumford-lem-compact-holomorphic-functions}
\dcref{M:ch1:I.1}
Every holomorphic map from a compact connected complex manifold to $\C$ is
constant.  Consequently a holomorphic map from a compact connected manifold
to a finite-dimensional complex vector space is constant.
\begin{proof}
The absolute value of each coordinate function attains a maximum.  The maximum
principle makes each coordinate locally, hence globally, constant on the
connected manifold.
\end{proof}
\notready
\end{lemma}

\begin{theorem}[Uniformization of compact complex Lie groups]
\label{mumford-thm-complex-lie-uniformization}
\dcref{M:ch1:I.1}
Let $X$ be a compact connected complex Lie group with identity $e$ and let
$V=\Lie_e(X)$.  The exponential homomorphism $\exp:V\to X$ is surjective,
its kernel $\Lambda$ is a lattice in $V$, and $\exp$ induces an isomorphism
of complex Lie groups $V/\Lambda\simeq X$.  In particular $X$ is commutative.
\uses{mumford-conv-complex-analytic,mumford-lem-compact-holomorphic-functions}
\notready
\end{theorem}

\begin{proposition}[Linear description of homomorphisms]
\label{mumford-prop-torus-homomorphisms}
\dcref{M:ch1:I.1}
For complex tori $X=V/\Lambda$ and $Y=W/\Gamma$, every holomorphic group
homomorphism $X\to Y$ has a unique complex-linear lift $u:V\to W$ with
$u(\Lambda)\subseteq\Gamma$.  Conversely every such $u$ descends to a
holomorphic group homomorphism.  Thus
\[
  \Hom(X,Y)=\{u\in\Hom_\C(V,W):u(\Lambda)\subseteq\Gamma\}.
\]
\uses{mumford-thm-complex-lie-uniformization}
\notready
\end{proposition}

\begin{corollary}[Topological cohomology of a torus]
\label{mumford-cor-torus-cohomology}
\dcref{M:ch1:I.1}
For $X=V/\Lambda$, the singular cohomology ring with integral coefficients is
canonically
\[
 H^\bullet(X,\Z)\simeq \bigwedge^\bullet\Hom_\Z(\Lambda,\Z),
\]
and the fundamental group is $\pi_1(X)\simeq\Lambda$.
\uses{mumford-lem-compact-lattice-quotient}
\notready
\end{corollary}

\section{Line bundles on a complex torus}
\label{mumford-I2}
\dcref{M:ch1:I.2}

\begin{definition}[Factor of automorphy]
\label{mumford-def-factor-automorphy}
\dcref{M:ch1:I.2}
Let $X=V/\Lambda$.  A factor of automorphy is a holomorphic function
$a:\Lambda\times V\to\C^\times$ satisfying
\[
 a(\lambda+\mu,z)=a(\lambda,z+\mu)a(\mu,z).
\]
It defines a line bundle $L_a$ by identifying $(z,t)$ with
$(z+\lambda,a(\lambda,z)t)$ on $V\times\C$.
\notready
\end{definition}

\begin{lemma}[Gauge equivalence]
\label{mumford-lem-factor-gauge}
\dcref{M:ch1:I.2}
Two factors $a$ and $a'$ define isomorphic line bundles if and only if there
is a nowhere-vanishing holomorphic $b:V\to\C^\times$ with
$a'(\lambda,z)=b(z+\lambda)a(\lambda,z)b(z)^{-1}$.
\uses{mumford-def-factor-automorphy}
\notready
\end{lemma}

\begin{definition}[Riemann form]
\label{mumford-def-riemann-form-analytic}
\dcref{M:ch1:I.2}
An alternating bilinear form $E:V_\R\times V_\R\to\R$ is a Riemann form for
$(V,\Lambda)$ if $E(\Lambda,\Lambda)\subseteq\Z$ and the real bilinear form
$H(u,v)=E(iu,v)+iE(u,v)$ is Hermitian positive definite.  Equivalently, $E$
is the imaginary part of a positive Hermitian form whose imaginary part is
integral on $\Lambda$.
\notready
\end{definition}

\begin{theorem}[Appell--Humbert classification]
\label{mumford-thm-appell-humbert}
\dcref{M:ch1:I.2}
For a complex torus $X=V/\Lambda$, every holomorphic line bundle is represented
by a pair $(H,\alpha)$, where $H$ is Hermitian and
$E=\operatorname{Im}H$ is integral on $\Lambda$.  In the factor convention
used here the second datum is a semicharacter:
\[
 \alpha(\lambda+\mu)=\alpha(\lambda)\alpha(\mu)
   \exp\!\bigl(\pi i\,E(\lambda,\mu)\bigr).
\]
The associated factor of automorphy is obtained from $H$ and $\alpha$ by the
standard Gaussian expression (with the same choice of which argument of $H$
is complex-linear).  Changing the trivialization gives exactly the gauge
relation in Lemma~$\ref{mumford-lem-factor-gauge}$.  The Chern class depends
only on $E$, and $L\mapsto E$
identifies $\NS(X)$ with the integral alternating forms of type $(1,1)$.
\uses{mumford-def-factor-automorphy,mumford-lem-factor-gauge,mumford-def-riemann-form-analytic}
\notready
\end{theorem}

\begin{proposition}[Chern class and translations]
\label{mumford-prop-analytic-c1-translation}
\dcref{M:ch1:I.2}
If $L$ corresponds to $(H,\alpha)$, then $c_1(L)$ is represented by the
constant form $E=\operatorname{Im}H$ on $V$ and translations act trivially on
$c_1(L)$.  The subgroup
\[
 K(L)=\{x\in X\mid t_x^*L\otimes L^{-1}\text{ is topologically trivial}\}
\]
is the kernel of the homomorphism induced by $E$ on $V/\Lambda$.
\uses{mumford-thm-appell-humbert}
\notready
\end{proposition}

\begin{theorem}[Theta-space index formula]
\label{mumford-thm-analytic-theta-index}
\dcref{M:ch1:I.2}
For a line bundle $L$ on a complex torus, the cohomology of $L$ is determined
by the signature of the associated Hermitian form.  If $E$ is nondegenerate,
then exactly one cohomology group $H^q(X,L)$ is nonzero, its index $q$ is the
number of negative directions of $H$, and
\[
 \sum_q(-1)^q\dim H^q(X,L)=\frac{c_1(L)^g}{g!}.
\]
For a positive form this gives $H^q(X,L)=0$ for $q>0$ and identifies the
dimension of the theta space with the degree of $L$.
\uses{mumford-thm-appell-humbert,mumford-prop-analytic-c1-translation}
\notready
\end{theorem}

\section{Algebraizability of tori}
\label{mumford-I3}
\dcref{M:ch1:I.3}

\begin{theorem}[Riemann's algebraizability criterion]
\label{mumford-thm-riemann-algebraizability}
\dcref{M:ch1:I.3}
A complex torus $X=V/\Lambda$ is the analytification of a projective algebraic
variety if and only if $(V,\Lambda)$ admits a Riemann form.  In that case a
line bundle whose Chern form is positive is ample.
\uses{mumford-def-riemann-form-analytic,mumford-thm-appell-humbert}
\notready
\end{theorem}

\begin{theorem}[Theta embedding]
\label{mumford-thm-theta-embedding}
\dcref{M:ch1:I.3}
If $L$ is a positive line bundle on a complex torus, then $L^{\otimes n}$ is
very ample for all sufficiently large $n$ and the morphism defined by a basis
of $H^0(X,L^{\otimes n})$ embeds $X$ in projective space.  The image is a
projective algebraic group variety.
\uses{mumford-thm-analytic-theta-index,mumford-thm-riemann-algebraizability}
\notready
\end{theorem}

\begin{corollary}[Complex abelian varieties]
\label{mumford-cor-complex-av-torus}
\dcref{M:ch1:I.3}
For a compact connected complex Lie group $X$, the following are equivalent:
\begin{enumerate}
\item $X$ is the analytification of an abelian variety;
\item $X$ is a polarizable complex torus;
\item $X$ has a positive line bundle.
\end{enumerate}
\uses{mumford-thm-complex-lie-uniformization,mumford-thm-riemann-algebraizability,mumford-thm-theta-embedding}
\notready
\end{corollary}

\begin{example}[Elliptic curves]
\label{mumford-ex-complex-elliptic}
\dcref{M:ch1:I.3}
For $\tau\in\mathfrak h$, the torus $\C/(\Z+\tau\Z)$ has the standard Riemann
form $E(1,\tau)=1$ and is the analytification of a genus-one projective curve.
Changing the lattice by $\operatorname{SL}_2(\Z)$ gives the usual moduli
identification of elliptic curves.
\uses{mumford-def-riemann-form-analytic,mumford-cor-complex-av-torus}
\notready
\end{example}

\section{Supporting dependency interfaces}
\label{mumford-I-coverage-register}
The following interfaces isolate auxiliary arguments used by the analytic
results in this chapter.  Their detailed proofs are treated as external inputs.

\begin{remark}[Complex-torus package]
\label{ch01-sec01-coverage}
The lattice quotient, its fundamental group, exterior-algebra cohomology, and
the linear lifting description of maps form the basic complex-torus package.
\dcref{M:ch1:I.1}
\uses{mumford-lem-compact-lattice-quotient,mumford-prop-torus-homomorphisms,
 mumford-cor-torus-cohomology}
\notready
\end{remark}

\begin{remark}[Uniformization proof interface]
\label{ch01-sec01-unread-results}
The exponential-kernel and period-lattice arguments behind uniformization are
retained as an external input.
\dcref{M:ch1:I.1}
\uses{mumford-thm-complex-lie-uniformization}
\notready
\end{remark}

\begin{remark}[Factors and integral forms]
\label{ch01-sec02-coverage}
Factors of automorphy, gauge changes, and integral alternating forms together
encode the analytic Picard group; the Appell--Humbert node fixes the semicharacter
convention used by the analytic comparison results.
\dcref{M:ch1:I.2}
\uses{mumford-def-factor-automorphy,mumford-lem-factor-gauge,
 mumford-thm-appell-humbert}
\notready
\end{remark}

\begin{remark}[Theta-dimension proof interface]
\label{ch01-sec02-unread-results}
The theta-dimension and positivity calculations depend on the factor data and
are recorded as external inputs.
\dcref{M:ch1:I.2}
\uses{mumford-thm-analytic-theta-index}
\notready
\end{remark}

\begin{remark}[Algebraizability package]
\label{ch01-sec03-coverage}
The Riemann criterion and theta embedding form the algebraizability package;
the elliptic example supplies a concrete one-dimensional instance.
\dcref{M:ch1:I.3}
\uses{mumford-thm-riemann-algebraizability,mumford-thm-theta-embedding}
\notready
\end{remark}

\begin{remark}[Algebraizability proof interface]
\label{ch01-sec03-unread-results}
Period-matrix relations, dimension computations, and the remaining examples are
external inputs rather than omitted dependencies.
\dcref{M:ch1:I.3}
\uses{mumford-thm-riemann-algebraizability,mumford-thm-theta-embedding}
\notready
\end{remark}

\begin{remark}[Lattice-torus package]
\label{ch01-sec01-lattice-torus}
The lattice quotient and compact-torus convention are represented by the
uniformization package.
\dcref{M:ch1:I.1}
\uses{mumford-conv-complex-analytic,mumford-lem-compact-lattice-quotient}
\notready
\end{remark}

\begin{remark}[Torus morphism package]
\label{ch01-sec01-morphisms}
The lift-to-vector-spaces description gives the complete homomorphism input
used by Tate theory and the analytic comparison results.
\dcref{M:ch1:I.1}
\uses{mumford-prop-torus-homomorphisms}
\notready
\end{remark}

\begin{remark}[Factor package]
\label{ch01-sec02-factors}
Factors of automorphy and their gauge equivalence are the analytic line-bundle
construction.
\dcref{M:ch1:I.2}
\uses{mumford-def-factor-automorphy,mumford-lem-factor-gauge}
\notready
\end{remark}

\begin{remark}[Riemann-form package]
\label{ch01-sec02-riemann-form}
The integral alternating form attached to a line bundle is the Chern-form datum
used by Appell--Humbert.
\dcref{M:ch1:I.2}
\uses{mumford-def-riemann-form-analytic,mumford-thm-appell-humbert}
\notready
\end{remark}

\begin{remark}[Algebraizability criterion package]
\label{ch01-sec03-algebraizability-criterion}
The existence of a positive integral Riemann form is equivalent to projective
algebraizability.
\dcref{M:ch1:I.3}
\uses{mumford-thm-riemann-algebraizability}
\notready
\end{remark}

\begin{remark}[Projective realization package]
\label{ch01-sec03-projective-realization}
A sufficiently high positive tensor power embeds the torus in projective space.
\dcref{M:ch1:I.3}
\uses{mumford-thm-theta-embedding}
\notready
\end{remark}
